%=============================================================================!
% 11.S  WHAT WE NOW HAVE                                                      !
%=============================================================================!
\unnumbered{What we now have}

A bond held at a fixed length adds a force along the old bond to each
step, and the bond at the end of the step is the provisional one plus a
multiple $\eta$ of the old: a quadratic in $\eta$ whose first-order
solution is a step of Newton's method. SHAKE corrects the bonds one at a
time with that step and sweeps until all hold, the error of a single bond
squaring at each sweep and that of a water molecule's triangle falling by
a steady factor.

RATTLE gives velocity Verlet the same constraints,
correcting the positions and half-step velocities after the drift and
removing each relative velocity's part along its bond after the second
half kick; both corrections keep the total momentum, the step is
reversible and symplectic, and it agrees with ASE to the rounding of the
arithmetic. Each independent held distance removes one freedom. On 64 molecules of
TIP3P water, at an energy fluctuation of 1\% of the kinetic energy's,
holding the O--H bonds lengthens the step from 0.41 to
\SI{0.97}{\femto\second} and holding all three distances to
\SI{1.98}{\femto\second}.

RESPA splits the forces and the Liouville
operator into fast and slow parts and applies the costly slow forces once
per outer step; it saves a factor of 1.8 in their calls for flexible
water, and its outer step is limited by resonance: for two springs the
outer step's matrix has the trace $2\cos\theta - (\omega_{\mathrm s}^2\dt
/\omega_{\mathrm f})\sin\theta$, unstable in bands just below multiples of
half the fast period, and in water beyond about a quarter of the fastest
period.

Moving mass from heavy atoms to their hydrogens slows the fastest
motions by a factor of about 1.6 to 1.7 while keeping the forces and, for
free or rigid molecules, every average over arrangements, and lengthens
the step of rigid water by a further 1.7. All
these gains are factors of a few: the Verlet-based methods used here
still need a step short enough to resolve the fastest motion they keep.

A healthy run shows $K$
and $U$ mirroring each other, a narrow band of total energy without drift,
bonds at the tolerance and a still centre of mass; eight common faults
each leave their own mark, caught by the checks of
\texttt{mdlab.diagnostics} and of the record.

Chapter~12 turns from single trajectories to the statistics of many
states: what temperature is, why the kinetic energy is shared equally
among the freedoms counted here, and what ensemble a simulation samples.
